%% ============================================================
%%  Artigo Científico — Disciplina de Pesquisa e Ordenação
%%  Universidade Tecnológica Federal do Paraná — UTFPR
%%  Formato: SBC (Sociedade Brasileira de Computação)
%% ============================================================

\documentclass[12pt]{article}

\usepackage{sbc-template}
\usepackage{graphicx,url}

\graphicspath{
  {../output/plots_python/}%  quando compilado de dentro de artigo-qmhr/
  {output/plots_python/}%     quando compilado da raiz do projeto
}

\usepackage[nil,shorthands=off]{babel}
\usepackage[utf8]{inputenc}
\usepackage{booktabs}
\usepackage{amsmath}
\usepackage{multirow}

\newcommand{\num}[1]{#1}

\sloppy

% ──────────────────────────────────────────────────────────────
%  Metadados
% ──────────────────────────────────────────────────────────────
\title{Análise Empírica de Algoritmos de Ordenação:\\
       Uma Avaliação Experimental de Quick Sort, Merge Sort,\\
       Heap Sort e Radix Sort em Entradas de Grande Volume}

\author{Gustavo R. Mazur\inst{1}}

\address{
  Departamento de Informática -- Universidade Tecnológica Federal do Paraná (UTFPR)\\
  \email{gustamomazur@alunos.utfpr.edu.br}
}

\begin{document}

\maketitle

% ──────────────────────────────────────────────────────────────
%  Resumo e Abstract
% ──────────────────────────────────────────────────────────────
\begin{abstract}
Sorting algorithms with sub-quadratic time complexity are fundamental
to modern computing, yet algorithms within the same asymptotic class
exhibit substantially different practical behavior due to memory access
patterns, cache utilization, and pivot or partitioning strategies.
This work presents an empirical study comparing four sorting algorithms
--- Quick Sort, Merge Sort, Heap Sort, and Radix Sort --- through a
reproducible benchmarking infrastructure implemented in Go~1.23.
Experiments were conducted on five input sizes (300,000 to 2,500,000
elements) under three dataset configurations (sorted, reversed, and
random), with three independent runs per combination. Metrics collected
include execution time, elementary operations (comparisons and element
movements), and heap memory allocations. The results confirm theoretical
complexity predictions and reveal relevant practical distinctions:
Radix Sort, a non-comparison-based algorithm, consistently outperforms
all others by circumventing the $\Omega(n \log n)$ lower bound; Quick Sort
with median-of-three pivot selection achieves its fastest execution on
sorted inputs; Heap Sort incurs the highest execution time despite
sharing the same asymptotic class as Quick Sort and Merge Sort due to
poor cache locality; and Merge Sort achieves maximum timing stability
at the cost of significant auxiliary memory allocation.

\noindent{\bf Keywords:} Sorting Algorithms, Empirical Analysis,
Computational Complexity, Quick Sort, Radix Sort, Go Language.
\end{abstract}

\begin{resumo}
Algoritmos de ordenação com complexidade de tempo sub-quadrática
são fundamentais para a computação moderna, porém algoritmos dentro
da mesma classe assintótica exibem comportamento prático substancialmente
distinto em razão de padrões de acesso à memória, localidade de
\textit{cache} e estratégias de particionamento ou seleção de pivô.
O presente trabalho apresenta um estudo empírico comparando quatro
algoritmos de ordenação --- Quick Sort, Merge Sort, Heap Sort e
Radix Sort --- por meio de uma infraestrutura de \textit{benchmarking}
reprodutível implementada em Go~1.23. Os experimentos foram conduzidos
sobre cinco tamanhos de entrada (300.000 a 2.500.000 elementos) sob
três configurações de \textit{dataset} (ordenado, invertido e aleatório),
com três execuções independentes por combinação. As métricas coletadas
incluem tempo de execução, operações elementares (comparações e
movimentações de elementos) e alocações de memória em \textit{heap}.
Os resultados confirmam as predições teóricas de complexidade e revelam
distinções práticas relevantes: o Radix Sort, único algoritmo
não-comparativo avaliado, supera consistentemente os demais ao
contornar o limite inferior $\Omega(n \log n)$; o Quick Sort com
seleção de pivô pela mediana de três atinge desempenho máximo em
entradas ordenadas; o Heap Sort apresenta o maior tempo de execução
apesar de compartilhar a mesma classe assintótica dos demais, em
razão de baixa localidade de \textit{cache}; e o Merge Sort alcança
máxima estabilidade temporal ao custo de alocações de memória auxiliar
significativas.

\noindent{\bf Palavras-chave:} Algoritmos de Ordenação, Análise Empírica,
Complexidade Computacional, Quick Sort, Radix Sort, Linguagem Go.
\end{resumo}

% ──────────────────────────────────────────────────────────────
%  1. INTRODUÇÃO
% ──────────────────────────────────────────────────────────────
\section{Introdução}

A escolha do algoritmo de ordenação adequado é uma decisão recorrente
em engenharia de software, impactando desde o desempenho de sistemas de
banco de dados até a eficiência de pipelines de processamento de dados
em larga escala \cite{cormen2009}. Embora a análise assintótica forneça
um primeiro critério de seleção, algoritmos pertencentes à mesma classe
de complexidade $O(n \log n)$ podem divergir substancialmente em
desempenho real em função de fatores que a notação $O(\cdot)$ abstrai
deliberadamente: localidade de \textit{cache}, custo de alocação de
memória, qualidade da seleção de pivô e sensibilidade à distribuição
de entrada \cite{knuth1998}.

Além dos algoritmos comparativos de classe $O(n \log n)$, existe uma
família distinta de algoritmos baseados em contagem ou distribuição —
como o Radix Sort — que contorna o limite inferior $\Omega(n \log n)$
provado para ordenação por comparação \cite{cormen2009}. Esses algoritmos
exploram propriedades estruturais dos dados (como a representação numérica
dos elementos) para atingir complexidade $O(nk)$, onde $k$ é o número
de dígitos da chave. Essa diferença fundamental requer avaliação empírica
para quantificar o ganho prático em cenários reais.

Diante desse contexto, o presente trabalho adota uma abordagem
experimental e reprodutível para comparar quatro algoritmos de
ordenação amplamente empregados. Especificamente, os objetivos são:

\begin{enumerate}
  \item Implementar e instrumentar os algoritmos Quick Sort, Merge Sort,
        Heap Sort e Radix Sort com contagem exata de comparações e
        movimentações de elementos;
  \item Executar os algoritmos sobre três tipos de \textit{dataset}
        (ordenado, invertido e aleatório) em cinco tamanhos de entrada
        (300.000 a 2.500.000 elementos);
  \item Comparar os resultados empíricos com as predições teóricas
        de complexidade;
  \item Identificar e discutir divergências práticas não capturadas
        pela análise assintótica, em particular a influência da
        distribuição de entrada e dos padrões de acesso à memória.
\end{enumerate}

Os algoritmos foram selecionados por representarem estratégias
distintas de ordenação: o Quick Sort baseia-se em particionamento
por pivô; o Merge Sort aplica a estratégia de divisão e conquista
com fusão; o Heap Sort utiliza uma estrutura de \textit{heap} binário;
e o Radix Sort emprega distribuição por dígito, dispensando
comparações entre elementos.

O restante deste artigo está organizado da seguinte forma:
a Seção~\ref{sec:theory} apresenta o referencial teórico dos
algoritmos avaliados; a Seção~\ref{sec:related} discute os
trabalhos relacionados; a Seção~\ref{sec:method} descreve a
metodologia experimental adotada; a Seção~\ref{sec:results}
apresenta os resultados obtidos; a Seção~\ref{sec:discussion}
analisa e discute esses resultados; e a Seção~\ref{sec:conclusion}
conclui o artigo.

% ──────────────────────────────────────────────────────────────
%  2. REFERENCIAL TEÓRICO
% ──────────────────────────────────────────────────────────────
\section{Referencial Teórico}
\label{sec:theory}

\subsection{Quick Sort}

O Quick Sort, proposto por Hoare em 1962 \cite{hoare1962}, é um
algoritmo de ordenação baseado em divisão e conquista que seleciona
um elemento \textit{pivô} e reorganiza o vetor de modo que todos os
elementos menores que o pivô precedam-no e todos os maiores
sucedam-no. O processo é aplicado recursivamente às duas subpartições
resultantes.

A complexidade de tempo do Quick Sort é $O(n \log n)$ no caso médio
e $O(n^2)$ no pior caso, que ocorre quando partições degeneradas são
geradas sistematicamente --- por exemplo, ao selecionar sempre o
primeiro ou último elemento como pivô em vetores já ordenados
\cite{cormen2009}. A seleção do pivô pela \textit{mediana de três}
--- o elemento mediano entre o primeiro, o central e o último do
subvetor --- mitiga eficazmente esse risco, garantindo partições
equilibradas mesmo em entradas ordenadas ou invertidas \cite{sedgewick2011}.
A complexidade de espaço é $O(\log n)$ em virtude da pilha de
recursão, não havendo alocação de memória auxiliar proporcional a $n$.

\subsection{Merge Sort}

O Merge Sort é um algoritmo de divisão e conquista que divide
recursivamente o vetor em dois sub-vetores de tamanhos aproximadamente
iguais, os ordena de forma independente e, em seguida, os funde
(\textit{merge}) em um único vetor ordenado \cite{cormen2009}.
A complexidade de tempo é $O(n \log n)$ para todos os casos (melhor,
médio e pior), uma vez que o número de níveis de recursão é
$\lceil \log_2 n \rceil$ e cada nível realiza $O(n)$ trabalho na
fusão \cite{knuth1998}.

Uma característica fundamental do Merge Sort é sua estabilidade:
elementos com chaves iguais mantêm a ordem relativa original do
vetor de entrada. A implementação de fusão requer um vetor auxiliar
de tamanho $O(n)$, o que implica complexidade de espaço $O(n)$
--- distinção relevante em sistemas com restrições de memória
\cite{goodrich2014}.

\subsection{Heap Sort}

O Heap Sort, introduzido por Williams em 1964 \cite{williams1964},
é um algoritmo de ordenação baseado em uma estrutura de dados de
\textit{heap} binário máximo (\textit{max-heap}). O algoritmo
opera em duas fases: primeiro, constrói o \textit{heap} a partir
do vetor de entrada (\textit{heapify}), em $O(n)$; depois,
extrai repetidamente o elemento máximo e o posiciona ao final
do vetor, realizando $O(\log n)$ operações de restauração do
\textit{heap} por extração.

A complexidade de tempo é $O(n \log n)$ para todos os casos, e
a complexidade de espaço é $O(1)$, pois o algoritmo opera
\textit{in-place} sem estruturas auxiliares adicionais
\cite{cormen2009}. Apesar dessas garantias teóricas, o Heap Sort
é reconhecidamente penalizado por baixa localidade de
\textit{cache} em hardware moderno: o acesso ao \textit{heap}
binário segue um padrão de saltos não sequenciais na memória
(de $i$ para $2i+1$ ou $2i+2$), que conflita com a hierarquia
de \textit{cache} e a lógica de \textit{prefetch} dos processadores
modernos \cite{biggar2008}.

\subsection{Radix Sort}

O Radix Sort é um algoritmo de ordenação baseado em distribuição
que não realiza comparações diretas entre elementos \cite{cormen2009}.
A variante LSD (\textit{Least Significant Digit}) processa o vetor
em $k$ passagens, onde $k$ é o número de dígitos da maior chave,
ordenando por dígito do menos para o mais significativo. Em cada
passagem, os elementos são distribuídos em baldes (\textit{buckets})
de acordo com o valor do dígito corrente e, em seguida, coletados
em ordem.

A complexidade de tempo é $O(nk)$, onde $n$ é o número de elementos
e $k$ é o número de dígitos na representação da chave. Quando $k$
é constante ou $O(\log n)$, o algoritmo supera os algoritmos
comparativos, que estão limitados inferiormente por $\Omega(n \log n)$
\cite{knuth1998}. A complexidade de espaço é $O(n + b)$, onde $b$
é a base numérica utilizada (tipicamente $b = 10$ ou $b = 256$),
em virtude da alocação dos vetores de contagem e do vetor de saída.
O Radix Sort não é diretamente aplicável a dados não-inteiros sem
transformação prévia.

% ──────────────────────────────────────────────────────────────
%  3. TRABALHOS RELACIONADOS
% ──────────────────────────────────────────────────────────────
\section{Trabalhos Relacionados}
\label{sec:related}

A literatura sobre análise experimental de algoritmos de ordenação
pode ser organizada em dois grupos principais: estudos empíricos
comparativos de desempenho e análises específicas da interação
entre algoritmos e características de hardware.

\subsection*{Estudos Empíricos Comparativos de Desempenho}

Pizarro-Vasquez et al.~\cite{pizarro2020} conduziram uma avaliação
empírica de onze algoritmos de ordenação --- incluindo Quick Sort,
Merge Sort, Heap Sort e Shell Sort --- implementados em C e Python.
Os experimentos utilizaram quatro configurações de entrada (aleatória,
quase ordenada, invertida e com poucos valores únicos) e avaliaram
tempos de execução. Os autores confirmaram que algoritmos de classe
$O(n \log n)$ superam consistentemente os de classe $O(n^2)$
à medida que o tamanho da entrada cresce. Contudo, o estudo não
instrumenta contagens exatas de operações elementares, não quantifica
a variabilidade estatística entre execuções independentes e não
avalia algoritmos não-comparativos como o Radix Sort.

Musser~\cite{musser1997} propôs o Introsort, algoritmo híbrido que
combina Quick Sort, Heap Sort e Insertion Sort para garantir
$O(n \log n)$ no pior caso com constante prática reduzida. A análise
demonstra que o Heap Sort, embora assintoticamente equivalente ao
Quick Sort, impõe sobrecarga significativa em cenários de pivô
degenerado. Esse trabalho evidencia que a diferença prática entre
algoritmos da mesma classe assintótica é suficientemente grande
para motivar soluções híbridas em implementações de produção ---
motivação direta para o estudo empírico conduzido no presente artigo.

\subsection*{Análise de Predição de Desvios e Localidade de Cache}

Biggar et al.~\cite{biggar2008} investigaram o comportamento de
algoritmos de ordenação clássicos sob a perspectiva de predição de
desvios em processadores modernos. Os autores demonstraram que
algoritmos com mesma complexidade assintótica podem divergir
substancialmente em desempenho prático em função do
\textit{miss rate} do preditor de desvios e da localidade de
\textit{cache}. Em particular, o estudo evidencia que o Heap Sort
apresenta o pior desempenho de \textit{cache} dentre os algoritmos
$O(n \log n)$, enquanto o Merge Sort, apesar de alocar memória
auxiliar, apresenta padrão de acesso sequencial favorável ao
\textit{prefetch} do processador.

Diferentemente dos trabalhos acima, o presente artigo adota
instrumentação \textit{in-code} que contabiliza exatamente as
comparações e movimentações realizadas por cada algoritmo,
incluindo o Radix Sort --- um algoritmo não-comparativo ausente
na maioria dos estudos experimentais identificados. A
infraestrutura de \textit{benchmarking} é implementada em Go
--- linguagem que compila para código de máquina nativo, sem
\textit{overhead} de JIT --- e avalia entradas de até
2.500.000 elementos com três execuções independentes por
combinação, permitindo quantificar a variabilidade das medições.

% ──────────────────────────────────────────────────────────────
%  4. METODOLOGIA
% ──────────────────────────────────────────────────────────────
\section{Metodologia}
\label{sec:method}

\subsection{Ambiente Experimental}

Todos os experimentos foram executados em hardware dedicado,
sem virtualização, conforme descrito na Tabela~\ref{tab:hardware}.
O sistema operacional não foi reiniciado entre as execuções,
mas o \textit{benchmarking} foi isolado por invocações explícitas
do Garbage Collector antes de cada medição.

\begin{table}[ht]
\centering
\caption{Especificações do ambiente de experimentação.}
\label{tab:hardware}
\begin{tabular}{ll}
\toprule
\textbf{Componente} & \textbf{Especificação} \\
\midrule
Modelo              & Dell G15 5515 \\
Processador         & AMD Ryzen 5 5600H, 3,3 GHz base, 6 núcleos / 12 threads \\
Arquitetura         & x86-64 \\
Memória RAM         & 16 GB DDR4 \\
Sistema Operacional & Windows 11 Home (Build 26200) \\
Linguagem           & Go 1.23 (compilação nativa, sem JIT) \\
\bottomrule
\end{tabular}
\end{table}

\subsection{Implementação e Instrumentação}
\label{subsec:impl}

A linguagem Go foi escolhida por ser compilada estaticamente para
código de máquina nativo (x86-64), eliminando o \textit{overhead}
de interpretação e o \textit{warm-up} de compiladores JIT,
presentes em linguagens como Java e C\#. O tempo medido reflete,
portanto, o custo computacional real das operações elementares.

Cada algoritmo foi implementado satisfazendo a interface
\texttt{SortAlgorithm}, que exige os métodos \texttt{Name() string}
e \texttt{Sort(arr []int) SortResult}. A estrutura \texttt{SortResult}
encapsula o vetor ordenado, o total de comparações e o total de
movimentações de elementos, contados incrementalmente durante a
execução --- instrumentação \textit{in-code} sem interceptação
externa.

O Quick Sort foi implementado com seleção de pivô pela mediana de
três elementos (primeiro, central e último do subvetor), garantindo
partições equilibradas em entradas ordenadas e invertidas. O Merge
Sort utiliza fusão com vetor auxiliar alocado por operação de fusão.
O Radix Sort utiliza a variante LSD (\textit{Least Significant
Digit}) em base decimal, com número de passagens determinado
dinamicamente pelo valor máximo do vetor de entrada. O Heap Sort
implementa o esquema clássico de \textit{heapify} descendente com
extração do máximo.

O motor de \textit{benchmarking} adota as seguintes garantias
de isolamento:

\begin{itemize}
  \item \textbf{Cópia limpa de dados:} cada execução recebe uma
        cópia independente do \textit{dataset} original, evitando
        que a ordenação in-place de uma execução afete as
        subsequentes.

  \item \textbf{Limpeza do Garbage Collector:} duas chamadas
        explícitas a \texttt{runtime.GC()} antes de cada medição
        eliminam o ruído de coletas de ciclos anteriores.

  \item \textbf{Medição de tempo de alta resolução:}
        \texttt{time.Since(start)} utiliza o relógio monotônico
        do sistema operacional com resolução de nanossegundos.

  \item \textbf{Medição de memória:} \texttt{runtime.MemStats}
        é capturado antes e após cada ordenação, contabilizando
        as alocações de \textit{heap} acumuladas durante a
        execução. O valor reportado representa alocações
        cumulativas, não o pico de memória em uso simultâneo.
\end{itemize}

A reprodutibilidade dos \textit{datasets} aleatórios é garantida
por semente fixa no gerador pseudoaleatório
(\texttt{rand.NewSource(42)}), assegurando que todos os experimentos
recebam exatamente os mesmos dados de entrada em qualquer execução.

\subsection{Design do Experimento}

Cada combinação (algoritmo $\times$ tipo de \textit{dataset} $\times$
tamanho de entrada) foi executada três vezes de forma independente.
Os tipos de \textit{dataset} e sua correspondência com os casos
teóricos são apresentados na Tabela~\ref{tab:datasets}.

\begin{table}[ht]
\centering
\caption{Tipos de \textit{dataset} e correspondência com casos teóricos.}
\label{tab:datasets}
\begin{tabular}{lll}
\toprule
\textbf{Dataset} & \textbf{Descrição} & \textbf{Caso teórico} \\
\midrule
\texttt{ordenado}  & $[1, 2, 3, \ldots, n]$         & Melhor caso (Quick Sort) \\
\texttt{invertido} & $[n, n{-}1, \ldots, 1]$        & Pior caso sem pivô mediana \\
\texttt{aleatório} & Permutação com seed fixa 42     & Caso médio \\
\bottomrule
\end{tabular}
\end{table}

Os tamanhos de entrada avaliados foram: 300.000, 450.000, 750.000,
1.000.000 e 2.500.000 elementos. Para cada grupo (algoritmo $\times$
\textit{dataset}), foram calculadas as seguintes estatísticas a
partir das três execuções: mínimo, máximo, média aritmética e
desvio padrão.

\subsection{Métricas Coletadas}

As métricas registradas por execução foram: (i) duração em
nanossegundos e milissegundos; (ii) número de comparações
elemento-a-elemento; (iii) número de movimentações de elementos
(trocas, deslocamentos ou distribuições para o Radix Sort); e
(iv) alocações de \textit{heap} em kilobytes. Os resultados foram
exportados para arquivos CSV e processados com a biblioteca
\textit{pandas} em Python para geração de gráficos
(\textit{matplotlib}/\textit{seaborn}) e tabelas estatísticas.

% ──────────────────────────────────────────────────────────────
%  5. RESULTADOS
% ──────────────────────────────────────────────────────────────
\section{Resultados}
\label{sec:results}

\subsection{Operações Elementares}

A Tabela~\ref{tab:ops} apresenta o número de comparações, movimentações
de elementos e alocações de memória para $n = 750.000$ elementos. Os
valores são determinísticos e idênticos nas três execuções de um mesmo
grupo. Para o Radix Sort, a coluna ``Movimentações'' contabiliza
distribuições de elementos entre baldes por passagem ($k \times n$),
não trocas entre posições do vetor --- distinção fundamental que
reflete a natureza não-comparativa do algoritmo.

\begin{table}[ht]
\centering
\caption{Operações elementares para $n = 750.000$ elementos.}
\label{tab:ops}
\resizebox{\textwidth}{!}{%
\begin{tabular}{llrrr}
\toprule
\textbf{Algoritmo} & \textbf{Dataset} &
\textbf{Comparações} & \textbf{Movimentações} & \textbf{Mem.\ (KB)} \\
\midrule
Quick Sort & ordenado  & \num{15.451.423} & \num{0}          & \num{0}       \\
Quick Sort & invertido & \num{15.451.424} & \num{375.000}    & \num{0}       \\
Quick Sort & aleatório & \num{20.874.735} & \num{3.396.991}  & \num{0}       \\
\midrule
Merge Sort & ordenado  & \num{7.180.608}  & \num{14.701.424} & \num{118.350} \\
Merge Sort & invertido & \num{7.520.816}  & \num{14.701.424} & \num{118.345} \\
Merge Sort & aleatório & \num{13.704.702} & \num{14.701.424} & \num{118.345} \\
\midrule
Heap Sort  & ordenado  & \num{27.735.697} & \num{14.563.298} & \num{0}       \\
Heap Sort  & invertido & \num{26.357.656} & \num{13.416.716} & \num{0}       \\
Heap Sort  & aleatório & \num{27.009.663} & \num{13.993.716} & \num{0}       \\
\midrule
Radix Sort & ordenado  & \num{0}          & \num{4.500.000}$^a$  & \num{73.501}  \\
Radix Sort & invertido & \num{0}          & \num{4.500.000}$^a$  & \num{73.501}  \\
Radix Sort & aleatório & \num{0}          & \num{5.250.000}$^a$  & \num{79.365}  \\
\bottomrule
\multicolumn{5}{l}{\small $^a$~Radix Sort LSD base 10: ordenado/invertido $=$ 6 dígitos
$\times$ 750.000; aleatório $=$ 7 dígitos $\times$ 750.000.}
\end{tabular}%
}
\end{table}

\subsection{Tempo de Execução}

A Tabela~\ref{tab:tempo} apresenta o tempo médio de execução em
milissegundos para todos os cinco tamanhos de entrada avaliados.
Os valores correspondem à média das três execuções independentes.

\begin{table}[ht]
\centering
\caption{Tempo médio de execução (ms) por algoritmo, \textit{dataset} e tamanho de entrada.}
\label{tab:tempo}
\resizebox{\textwidth}{!}{%
\begin{tabular}{llrrrrr}
\toprule
\textbf{Algoritmo} & \textbf{Dataset} &
\textbf{300k} & \textbf{450k} & \textbf{750k} &
\textbf{1.000k} & \textbf{2.500k} \\
\midrule
Quick Sort & ordenado  & 4,407  & 7,543  & 11,620  & 15,799  & 42,589  \\
Quick Sort & invertido & 4,945  & 7,573  & 11,162  & 16,401  & 42,730  \\
Quick Sort & aleatório & 21,459 & 32,467 & 56,313  & 75,809  & 200,900 \\
\midrule
Merge Sort & ordenado  & 17,584 & 25,181 & 42,335  & 55,819  & 157,627 \\
Merge Sort & invertido & 15,798 & 23,732 & 38,298  & 52,705  & 145,406 \\
Merge Sort & aleatório & 33,376 & 50,503 & 83,298  & 112,393 & 297,669 \\
\midrule
Heap Sort  & ordenado  & 27,942 & 28,583 & 44,543  & 61,510  & 164,793 \\
Heap Sort  & invertido & 20,198 & 31,801 & 64,463  & 70,042  & 167,258 \\
Heap Sort  & aleatório & 33,150 & 50,025 & 92,327  & 126,829 & 430,186 \\
\midrule
Radix Sort & ordenado  & 12,434 & 16,402 & 27,519  & 43,138  & 104,497 \\
Radix Sort & invertido & 11,612 & 17,082 & 27,197  & 40,540  & 99,828  \\
Radix Sort & aleatório & 14,403 & 19,450 & 31,063  & 47,279  & 111,373 \\
\bottomrule
\end{tabular}%
}
\end{table}

As Figuras~\ref{fig:300k}, \ref{fig:750k} e~\ref{fig:2500k} apresentam
os tempos médios de execução em escala logarítmica (base 10) para
entradas de 300.000, 750.000 e 2.500.000 elementos, respectivamente.
A escala logarítmica é necessária para a visualização simultânea de
algoritmos cujos tempos diferem em até uma ordem de grandeza para
o mesmo tamanho de entrada.

\begin{figure}[ht]
\centering
\includegraphics[width=0.82\textwidth]{300000/py_tempo_medio.png}
\caption{Tempo médio de execução para $n = 300.000$ elementos
         (escala $\log_{10}$).}
\label{fig:300k}
\end{figure}

\begin{figure}[ht]
\centering
\includegraphics[width=0.82\textwidth]{750000/py_tempo_medio.png}
\caption{Tempo médio de execução para $n = 750.000$ elementos
         (escala $\log_{10}$).}
\label{fig:750k}
\end{figure}

\begin{figure}[ht]
\centering
\includegraphics[width=0.82\textwidth]{2500000/py_tempo_medio.png}
\caption{Tempo médio de execução para $n = 2.500.000$ elementos
         (escala $\log_{10}$).}
\label{fig:2500k}
\end{figure}

% ──────────────────────────────────────────────────────────────
%  6. ANÁLISE E DISCUSSÃO
% ──────────────────────────────────────────────────────────────
\section{Análise e Discussão}
\label{sec:discussion}

\subsection{Confirmação Empírica das Complexidades Teóricas}

Os resultados obtidos confirmam as predições teóricas de complexidade
para os três algoritmos comparativos. Para o Merge Sort com entrada
aleatória e $n = 750.000$, o número de comparações observado foi
\num{13.704.702}, valor próximo a $n \log_2 n \approx 14.640.000$ ---
razão de $0{,}94$, consistente com o comportamento esperado de
$\Theta(n \log n)$ \cite{knuth1998}.

O Heap Sort apresentou contagens de comparações praticamente
constantes entre os três tipos de \textit{dataset}: \num{27.735.697}
(ordenado), \num{26.357.656} (invertido) e \num{27.009.663} (aleatório)
para $n = 750.000$. Esse comportamento decorre da invariante do
algoritmo: as operações de \textit{heapify} são determinadas
estruturalmente pelo tamanho do \textit{heap}, independentemente
da distribuição inicial dos elementos. O valor observado de
aproximadamente $1{,}84 \cdot n \log_2 n$ reflete o custo de duas
comparações por nível na restauração do \textit{heap}.

Para o Quick Sort com seleção de pivô pela mediana de três, as
comparações em entrada ordenada foram \num{15.451.423} ($\approx
1{,}05 \cdot n \log_2 n$), evidenciando que a estratégia de pivô
garante comportamento $O(n \log n)$ mesmo no tipo de entrada que
levaria o Quick Sort ingênuo ao pior caso.

O Radix Sort realizou zero comparações em todos os cenários,
confirmando empiricamente sua natureza não-comparativa.

\subsection{Quick Sort: Inversão de Desempenho por Seleção de Pivô}

Um dos achados mais relevantes deste estudo é o comportamento
contra-intuitivo do Quick Sort em relação à distribuição de entrada.
Para $n = 750.000$, o tempo médio em entrada \textit{ordenada}
foi de apenas 11,620\,ms, enquanto o tempo médio em entrada
\textit{aleatória} foi de 56,313\,ms --- uma diferença de fator
$4{,}85\times$, com a entrada ordenada sendo significativamente
\textit{mais rápida}.

Esse comportamento é explicado pela interação entre a estratégia
de mediana de três e a estrutura do vetor de entrada. Em um vetor
completamente ordenado, o elemento mediano entre o primeiro, o
central e o último é sempre exatamente a mediana do subvetor,
produzindo partições perfeitamente balanceadas em cada nível de
recursão. O resultado é o menor número possível de comparações
por nível e ausência de trocas (0 trocas para entrada ordenada
versus \num{3.396.991} para aleatória). Em entrada aleatória,
a mediana de três é uma boa aproximação, mas não garante
particionamento perfeito, resultando em maior variabilidade e
número total de operações.

O mesmo padrão é observado em entrada \textit{invertida}: 11,162\,ms
para $n = 750.000$, também substancialmente mais rápido que o caso
aleatório. Para $n = 2.500.000$, a diferença persiste: 42,589\,ms
(ordenado) e 42,730\,ms (invertido) contra 200,900\,ms (aleatório)
--- fator $4{,}71\times$. Esses resultados demonstram que, com
seleção adequada de pivô, o Quick Sort atinge seu melhor desempenho
prático justamente nos tipos de entrada que representam o pior
caso da variante ingênua.

\subsection{Radix Sort: Ruptura com o Limite Inferior Comparativo}

O Radix Sort foi o algoritmo mais rápido em todos os tamanhos
e tipos de \textit{dataset} avaliados. Para $n = 750.000$ em
entrada aleatória, o tempo médio foi de 31,063\,ms, contra
56,313\,ms do Quick Sort (fator $1{,}81\times$). Para
$n = 2.500.000$, o Radix Sort executou em 111,373\,ms contra
200,900\,ms do Quick Sort (fator $1{,}80\times$), indicando
que a vantagem relativa é estável com o crescimento de $n$.

A origem dessa vantagem é estrutural: o Radix Sort não está
sujeito ao limite inferior $\Omega(n \log n)$ imposto à ordenação
por comparação \cite{cormen2009}. Com $k = 6$ passagens para
entradas ordenadas e invertidas ($n \leq 10^6$, seis dígitos
decimais) e $k = 7$ passagens para entradas aleatórias (valores
aleatórios atingindo sete dígitos), a complexidade efetiva é
$O(7n)$, linear no tamanho da entrada.

Uma consequência notável é a insensibilidade à distribuição de entrada:
os tempos para ordenado (27,519\,ms), invertido (27,197\,ms) e
aleatório (31,063\,ms) diferem apenas em razão do número de dígitos
do valor máximo (seis versus sete passagens). Em contraste, os
algoritmos comparativos apresentam variações de até $4{,}85\times$
entre os tipos de \textit{dataset}.

O custo do Radix Sort é a memória auxiliar: \num{79.365}\,KB
($\approx 77$\,MB) para $n = 750.000$ em entrada aleatória, e
\num{256.493}\,KB ($\approx 250$\,MB) para $n = 2.500.000$.
Esse custo é significativo em sistemas embarcados ou ambientes com
memória restrita, tornando o algoritmo inadequado em tais contextos
mesmo sendo mais rápido.

\subsection{Heap Sort: Custo Oculto da Desordem de Memória}

O Heap Sort compartilha a classe de complexidade $O(n \log n)$ com
Quick Sort e Merge Sort, porém apresentou consistentemente os maiores
tempos de execução entre os algoritmos comparativos. Para
$n = 750.000$ em entrada aleatória, o Heap Sort consumiu 92,327\,ms,
contra 83,298\,ms do Merge Sort e 56,313\,ms do Quick Sort ---
fatores $1{,}11\times$ e $1{,}64\times$, respectivamente. Para
$n = 2.500.000$, o gap se amplia: 430,186\,ms (Heap Sort) contra
200,900\,ms (Quick Sort), fator $2{,}14\times$.

Essa degradação relativa está associada ao padrão de acesso à
memória inerente à estrutura do \textit{heap} binário. Cada operação
de \textit{sift-down} percorre um caminho da raiz para uma folha,
acessando posições $i$, $2i+1$ e $2i+2$ do vetor --- um padrão de
saltos que não favorece a hierarquia de \textit{cache} dos
processadores modernos \cite{biggar2008}.

A alta variabilidade observada em entrada invertida --- desvio padrão
de 17,320\,ms para uma média de 64,463\,ms (27\%) em $n = 750.000$
--- é consistente com padrões irregulares de \textit{cache miss}
induzidos pela dinâmica específica do \textit{heapify} sobre esse
tipo de entrada. O Heap Sort é, em contrapartida, o único algoritmo
avaliado que combina complexidade $O(n \log n)$ garantida no pior
caso com complexidade de espaço $O(1)$ --- propriedade relevante
em contextos onde memória é escassa e garantias de pior caso são
essenciais.

\subsection{Merge Sort: Estabilidade Temporal e Custo de Memória}

O Merge Sort apresentou a menor variabilidade temporal entre todos
os algoritmos avaliados. Para $n = 750.000$ em entrada aleatória,
o desvio padrão foi de apenas 0,281\,ms ($0{,}34\%$ da média de
83,298\,ms), evidenciando comportamento altamente previsível ---
característica valiosa em sistemas de tempo real.

No entanto, a previsibilidade vem acompanhada de um custo de memória
significativo. As alocações cumulativas de \textit{heap} atingiram
\num{118.350}\,KB ($\approx 116$\,MB) para $n = 750.000$ e
\num{433.784}\,KB ($\approx 424$\,MB) para $n = 2.500.000$ ---
valores que podem ser proibitivos em sistemas embarcados. O padrão
de alocações cumulativas, em lugar do pico de uso simultâneo,
amplifica os valores reportados: cada operação de fusão aloca e
libera um vetor auxiliar de tamanho proporcional ao subvetor
processado, resultando em alocações totais de $O(n \log n)$.

O Merge Sort também apresentou o tempo médio mais baixo em entradas
\textit{invertidas} entre os algoritmos $O(n \log n)$ comparativos:
38,298\,ms em $n = 750.000$, inferior ao resultado para entrada
\textit{ordenada} (42,335\,ms). Essa inversão é atribuível à
dinâmica das operações de fusão: para entradas invertidas, os
subvetores de cada nível de recursão são sequências decrescentes
perfeitamente regulares, permitindo fusões com padrão de acesso
favorável ao \textit{prefetch}.

\subsection{Estabilidade das Medições}

O desvio padrão das medições oferece uma estimativa da variabilidade
introduzida pelo ambiente de execução. O algoritmo com maior
instabilidade absoluta foi o Heap Sort em $n = 2.500.000$ com
entrada aleatória (desvio padrão de 20,968\,ms, ou 4,9\% da média
de 430,186\,ms). O Merge Sort, em contrapartida, atingiu desvio
padrão de apenas 0,140\,ms para $n = 1.000.000$ em entrada aleatória
(0,12\% da média).

O Quick Sort apresentou desvio padrão inferior a 1\% em todos
os cenários de entrada ordenada e invertida, e valores entre
0,4\% e 1,6\% para aleatória, confirmando comportamento estável
com a estratégia de mediana de três. O Radix Sort manteve desvio
padrão abaixo de 3\% em todos os cenários, consistente com seu
comportamento determinístico por passagem.

% ──────────────────────────────────────────────────────────────
%  7. CONCLUSÃO
% ──────────────────────────────────────────────────────────────
\section{Conclusão}
\label{sec:conclusion}

Este trabalho apresentou uma avaliação empírica e reprodutível de
quatro algoritmos de ordenação sobre entradas de até 2.500.000
elementos. Os experimentos confirmaram quantitativamente as predições
da análise assintótica e revelaram distinções práticas relevantes
que a notação $O(\cdot)$ não captura.

Os principais achados podem ser sintetizados em quatro constatações.
Em primeiro lugar, o Radix Sort supera consistentemente todos os
algoritmos comparativos por contornar o limite inferior $\Omega(n \log n)$
--- com fator $1{,}8\times$ em relação ao Quick Sort para $n = 2.500.000$
aleatório --- ao custo de memória auxiliar significativa ($\approx 250$\,MB
para 2.500.000 elementos). Em segundo lugar, a estratégia de pivô
determina o comportamento prático do Quick Sort de forma mais
significativa que o tipo de entrada: com mediana de três, entradas
ordenadas e invertidas são processadas $4{,}7\times$ mais rápido
que entradas aleatórias. Em terceiro lugar, algoritmos de mesma
classe $O(n \log n)$ podem divergir em até $2{,}1\times$ (Heap Sort
versus Quick Sort em $n = 2.500.000$), com o Heap Sort penalizado
pela baixa localidade de \textit{cache} apesar de ser o único
algoritmo in-place com pior caso garantido. Em quarto lugar, o Merge
Sort entrega a maior estabilidade temporal (desvio padrão de 0,14\,ms
em $n = 1.000.000$), mas impõe o maior custo de memória cumulativa
entre os algoritmos avaliados.

Este trabalho apresenta algumas limitações. Os experimentos foram
conduzidos em um único hardware com arquitetura x86-64; comportamentos
de \textit{cache} e \textit{pipeline} podem diferir em processadores
ARM ou em ambientes com hierarquias de memória distintas. O Radix Sort
foi avaliado apenas para inteiros positivos de até sete dígitos
decimais; sua extensão a inteiros negativos ou a chaves de comprimento
variável requereria adaptação da implementação. As alocações de
memória reportadas são cumulativas, não representando o pico de
uso simultâneo --- métrica complementar relevante para análise
de pressão de memória.

Como trabalhos futuros, pretende-se avaliar algoritmos híbridos de
produção como Introsort e Timsort, que combinam Quick Sort, Heap Sort
e Insertion Sort para garantir $O(n \log n)$ no pior caso com
constante prática mínima. Adicionalmente, planeja-se estender a
análise ao Counting Sort e ao Bucket Sort para dados de distribuição
uniforme, e conduzir testes estatísticos formais de comparação de
médias (Wilcoxon signed-rank) para validar a significância estatística
das diferenças observadas.

% ──────────────────────────────────────────────────────────────
%  REFERÊNCIAS
% ──────────────────────────────────────────────────────────────
\bibliographystyle{sbc}
\bibliography{referencias}

\end{document}
